\chapter{Resultados preditivos}\label{cap:resultados}
\section{Comparação global das quatro configurações}
A avaliação foi concluída para 20 partes Ridge e dez referências, com uma previsão externa por faixa e por variante. A auditoria posterior reproduziu as métricas e reconciliou as partes com o arquivo global. A Tabela~\ref{tab:metricas} apresenta os resultados da comparação principal e a sensibilidade de avaliação ponderada.

\begin{table}[htbp]
\centering\small
\caption{Desempenho nas previsões externas reunidas ($n=77.535$).}
\label{tab:metricas}
\begin{tabular}{lrrrr}
\toprule
Configuração & MAE & RMSE & R² & MAE ponderado \\
\midrule
Baseline & 8,3525 & 10,8255 & 0,5087 & 8,2194 \\
Tradicional & 8,2957 & 10,7538 & 0,5152 & 8,1652 \\
Profundo & 8,1285 & 10,5540 & 0,5330 & 8,0032 \\
Combinado & 8,1151 & 10,5381 & 0,5344 & 7,9908 \\
Média do treino & 12,5621 & 15,4463 & -0,0003 & 12,4181 \\
Mediana do treino & 12,4786 & 15,5708 & -0,0165 & 12,3213 \\
\bottomrule
\end{tabular}
\fonte{Elaboração própria a partir dos agregados auditados do estudo.}
\nota{Pesos iguais no ajuste e na avaliação principal. Última coluna: sensibilidade com pesos originais nas mesmas previsões; não houve retreino ponderado.}
\end{table}


O Combinado apresentou o menor MAE observado, 8,1151, seguido do Profundo, 8,1285. Tradicional e Baseline tiveram MAE de 8,2957 e 8,3525. O Combinado também apresentou os menores RMSE e maior $R^2$ entre as quatro configurações, mas essas métricas complementares não foram usadas para escolher retrospectivamente o critério de destaque. As referências de média e mediana do treino tiveram erros maiores.

Em relação ao Baseline, a diferença de MAE do Combinado foi de $-0{,}23736$ ponto, equivalente a redução relativa de 2,84\%. É um incremento médio de magnitude moderada na escala de zero a cem; não significa que a letra explique causalmente essa parcela da popularidade nem que todo caso tenha melhorado. O $R^2$ externo descreve erro relativo à variabilidade observada, sem interpretação como fração causal atribuível às letras.

\figura{comparacao_global}{Desempenho médio e incrementos nas previsões externas.}{fig:global}{Agregados da comparação Ridge atual. Pesos originais são usados somente como sensibilidade na avaliação.}
\FloatBarrier

\section{Incrementos específicos}
A Tabela~\ref{tab:contrastes} discrimina os cinco contrastes. O Profundo reduz o MAE frente ao Baseline em cerca de 0,2240 ponto, enquanto o Tradicional reduz cerca de 0,0568. O Combinado acrescenta apenas 0,0134 ponto de redução em relação ao Profundo. Assim, a ordenação pontual favorece a união das representações, mas quase todo o ganho conjunto frente ao Baseline já aparece ao acrescentar o embedding.

\begin{table}[htbp]
\centering\small
\caption{Diferenças pareadas de MAE entre configurações.}
\label{tab:contrastes}
\begin{tabular}{lrr}
\toprule
Contraste & Pesos iguais & Pesos originais \\
\midrule
Tradicional - Baseline & -0,05679 & -0,05420 \\
Profundo - Baseline & -0,22396 & -0,21627 \\
Combinado - Baseline & -0,23736 & -0,22861 \\
Combinado - Tradicional & -0,18057 & -0,17441 \\
Combinado - Profundo & -0,01340 & -0,01234 \\
\bottomrule
\end{tabular}
\fonte{Elaboração própria a partir dos agregados auditados do estudo.}
\nota{Diferença negativa favorece a primeira configuração; são contrastes observados, sem inferência integral de treinamento.}
\end{table}


Esse resultado é compatível com sobreposição informacional entre descritores e embeddings, mas não a demonstra como mecanismo. A família Ridge, suas escalas, a grade e a representação efetivamente usada participam do resultado. Não é possível concluir que descritores explícitos sejam sempre redundantes ou que embeddings sejam universalmente superiores em outras populações ou algoritmos.

Com pesos originais, o MAE do Combinado foi 7,9908 frente a 8,2194 do Baseline, redução relativa de 2,78\%. A ordenação é mantida, e o incremento Combinado--Profundo continua pequeno. Essa concordância é evidência de sensibilidade limitada à reponderação específica dos casos observados; não valida correção do viés de disponibilidade textual.

\section{Variação entre avaliações externas}
O Combinado apresentou menor MAE nos cinco folds externos, mantendo a mesma ordenação global. As diferenças frente ao Baseline variaram de aproximadamente $-0{,}2724$ a $-0{,}2070$. A Tabela~\ref{tab:folds} e a Figura~\ref{fig:folds} permitem examinar essa consistência sem tratar folds como observações independentes para um teste de hipótese.

\begin{table}[htbp]
\centering\small
\caption{MAE por avaliação externa nos mesmos componentes.}
\label{tab:folds}
\begin{tabular}{lrrrr}
\toprule
Fold & Baseline & Tradicional & Profundo & Combinado \\
\midrule
1 & 8,3287 & 8,2526 & 8,1155 & 8,0991 \\
2 & 8,2666 & 8,1989 & 8,0549 & 8,0392 \\
3 & 8,4241 & 8,3673 & 8,2364 & 8,2171 \\
4 & 8,3633 & 8,3211 & 8,1188 & 8,1128 \\
5 & 8,3798 & 8,3385 & 8,1171 & 8,1074 \\
\bottomrule
\end{tabular}
\fonte{Elaboração própria a partir dos agregados auditados do estudo.}
\nota{A numeração 1--5 corresponde a 0--4 nos arquivos. Folds não são réplicas independentes.}
\end{table}

\figura{estabilidade_folds}{MAE nas cinco avaliações externas agrupadas.}{fig:folds}{Mesmas observações e partições para as quatro configurações; a sobreposição dos treinos impede interpretar os cinco pontos como réplicas independentes.}
\FloatBarrier

\section{Penalização selecionada e limites do procedimento}
A Tabela~\ref{tab:alphas} registra os parâmetros escolhidos nas validações internas. Seis seleções ocorreram no extremo inferior da grade, 0,1. Esse fato é documentado como limitação da busca fixada e não autoriza expandi-la depois de observar os testes externos. A comparação responde pelo procedimento efetivamente executado, e não pelo melhor resultado que poderia ser obtido após novas tentativas.

\begin{table}[htbp]
\centering\small
\caption{Parâmetro de penalização selecionado internamente.}
\label{tab:alphas}
\begin{tabular}{lrrrr}
\toprule
Fold & Baseline & Tradicional & Profundo & Combinado \\
\midrule
1 & 1,0 & 1,0 & 1,0 & 1,0 \\
2 & 1,0 & 1,0 & 0,1 & 1,0 \\
3 & 0,1 & 0,1 & 0,1 & 0,1 \\
4 & 1,0 & 1,0 & 0,1 & 1,0 \\
5 & 1,0 & 10,0 & 1,0 & 1,0 \\
\bottomrule
\end{tabular}
\fonte{Elaboração própria a partir dos agregados auditados do estudo.}
\nota{A grade foi fixada antes dos resultados externos; não foi expandida após observar escolhas no extremo.}
\end{table}


Não foi ajustado automaticamente um novo modelo final em toda a coorte. Os resultados reportados são de avaliação externa agrupada, com os próprios tamanhos de treino definidos pela CV. Um modelo de implantação, se fosse objetivo posterior, teria outra finalidade e não herdaria sem ressalvas a mesma quantificação de erro ou incerteza.
